\subsubsection{Encodings}
Typical text encodings are \texttt{CSV}, \texttt{XML}\footnote{you have to be quite a bit of a moron to think that this is the best of the text encodings, but all right}
and \texttt{JSON}. Of course, there are also the (in my objective opinion) best formats\footnote{best for configuration files, not (necessarily) in general}
like \texttt{conf} and \texttt{yaml}, as well as \texttt{TOML}, \texttt{INI} and more.

These are good for being human readable, but not if you have terabytes of data.
From a practical example, I know that a 1GB JSON file in Python's built in parser takes about 1s to load.

This is why we need (direct) binary encodings for the data, such as the Apache formats.

Examples include Avro (row-oriented format, including a schema describing the fields and types, this is the header of the file), Parquet (Compressed column store),
CarbonData (Compressed column store with indexes) and more.

When designing formats, decisions need to be made on the file metadata, format layout, type system, encoding schemes, compression, filters and how to nest the data.

\bi{Metadata} Ideally the files are self-contained to allow for portability, and for space efficiency reasons (even though it may be beneficial for other reasons),
schemas are only stored once per file, in either the header (e.g. for Avro) or in the footer (e.g. Parquet). Other metadata however is more oftne stored once per frame.

The \bi{format layout} is the physical organization of the bytes on the disk. We can have row stores, column stores or hybrids:
\begin{itemize}
    \item N-ary Storage Model (NSM): All attributes stored contiguously in a single page (Row store)
    \item Decomposition Storage Model (DSM): Store single attributes for all tuples contiguously in a block of data (Column store)
    \item PAX: Horizontally partition data into row groups, then vertically partition the attributes into column chunks.
\end{itemize}

For the \bi{type system}, we choose what kind of types we want to support, be it primitives such as numbers and strings, or more complex types such as arrays.
Furthermore, we choose how they are physically represented, such as IEEE-754 encoding for floating point numbers, and logically, if they are complex types based on other types,
i.e. how they abstract a more primitive type.

For \bi{encoding schemes} include:
\begin{itemize}
    \item \textit{Dictionary encoding}: We find all unique items of a column and build a dictionary with integer keys and then only store the integers.
          This is ideal for columns with few distinct values and worse than not doing anything for columns with the UNIQUE constraint set.
    \item \textit{Run-Length encoding}: If we have many repeating values, called runs, we replace each run with a pair \texttt{(count, value)}.
          This can be done on the data itself, or on the binary representation of said data.
    \item \textit{Delta Encoding}: Store the first value (as-is), then store the difference between this value and the next.
          Since these deltas are often small, we can fit them into fewer bits.
\end{itemize}

For \bi{Block compression}, we can use typical compression algorithms, such as \texttt{LZ4} (which by the way is VERY impressive for compression ratio),
\texttt{Snappy}, \texttt{Zstd} (which \texttt{pacman} uses for the packages (\texttt{.pkg.tar.zst} files)), or older algorithms such as \texttt{gzip} or \texttt{xz}
(if you remember the whole malware debacle they had, you may not want to use that anymore).


\bi{Filters} in the files are primarily used to improve traverse performance, by filtering out pages that simply don't need to be searched, using some metadata.
For that, statistics are stored in the metadata, such as the min/max, null count, cardinality and more.
